% One shield, one rule and possibly one wire.
\begin{tikzpicture}[font=\sffamily\small, x=1cm, y=1cm]

% ---------------- the board and the fitted shield
\node[board, minimum width=50mm, minimum height=30mm] (nucleo) at (0,0) {};
\node[text=white, font=\sffamily\bfseries\small] at (0,-1.05) {NUCLEO-H7A3ZI-Q};
\node[hat, minimum width=44mm, minimum height=17mm] (shield) at (0,0.55) {};
\node[text=white, font=\sffamily\bfseries\small] at (0,0.95) {inertial shield};
\node[text=white, font=\sffamily\tiny] at (0,0.55) {two-wire bus, plus an interrupt line};
\node[pinrow, minimum width=40mm] at (0,1.35) {};
\node[pinrow, minimum width=40mm] at (0,-0.28) {};

% the data-ready line
\draw[cable, draw=red!60!black] (1.6,0.2) -- ++(1.5,0) -- ++(0,-0.9) -- ++(-1.5,0);
\node[font=\sffamily\tiny, anchor=west] at (3.2,-0.25) {data-ready};

% ---------------- the rule
\node[absentblk, text width=44mm] (rule) at (6.6,1.4)
  {never stack the two shields\\{\scriptsize they share addresses. The bus then
   answers for the wrong part, and the symptom is plausible data from a sensor
   that is not fitted}};

% ---------------- the three routing cases
\node[chlabel, anchor=north west] at (-4.6,-1.9) {getting the data-ready line onto a capture channel};
\node[pinused/.try, draw=linecol, fill=usrcol, text width=40mm, align=center,
      minimum height=8mm, font=\sffamily\scriptsize, anchor=north west] (c1) at (-4.6,-2.3)
  {preferred: it already lands on a capture-capable pin};
\node[draw=linecol, fill=hwcol, text width=40mm, align=center, minimum height=8mm,
      font=\sffamily\scriptsize, anchor=north west] (c2) at (-0.1,-2.3)
  {fallback 1: one jumper wire to a pin that can};
\node[draw=linecol, fill=black!6, text width=40mm, align=center, minimum height=8mm,
      font=\sffamily\scriptsize, anchor=north west] (c3) at (4.4,-2.3)
  {fallback 2: an external interrupt, and report the added latency};

\node[umlnote, text width=58mm, anchor=north west] at (-4.6,-3.6)
  {Which of the three applies is a board question rather than a shield question,
   and it is on chapter 4's confirm list: the shield's manual says which of its
   pins carries the interrupt, and the board manual says what that header
   position can be routed to. The chapter works in all three cases and reports
   which one was used, because the third one changes the timestamp error and the
   reader is entitled to know which number they are looking at.};
\node[umlnote, text width=44mm, anchor=north west] at (3.4,-3.6)
  {This is the only wiring in the chapter. Everything else is the shield sitting
   on the header the way it was designed to.};
\end{tikzpicture}
